% GENERATED FILE. Edit canonical lessons, then run npm run generate:books.
% canonical-source-hash: fnv1a64:0a04834dbc5c5980
% canonical-lessons: UR-C418-hisab, UR-C418-sud, UR-C418-qist, UR-C418-adaigi, UR-C418-ghatana

\chapter{Accounts and Instalments}
\label{ch:ur-accounts-and-instalments}
\hlchaptermodality{418}

\begin{chapteropening}
\textbf{By the end of this chapter:} \emph{I can talk about keeping accounts, paying interest and paying in instalments.}

Complete the last lesson of chapter 418: I can talk about keeping accounts, paying interest and paying in instalments.
\end{chapteropening}

\section[\texorpdfstring{\ur{حساب} (hisāb)}{حساب (hisāb)}]{\ur{حساب} (hisāb) — an account; a reckoning}

\label{lesson:UR-C418-hisab}



\begin{quote}
Before the new one: say the Urdu for a cheque, then the Urdu for to withdraw (money).
\end{quote}

\subsection*{You'll want to know: \ur{حساب}}
\textbf{\ur{حساب}} — \emph{hisāb} — \textquotedblleft{}an account; a reckoning\textquotedblright{}.

\begin{quote}
\ur{دکاندار} \ur{ہر} \ur{شام} \ur{اپنا} \ur{حساب} \ur{لکھتا} \ur{ہے۔}
\end{quote}

\begin{quote}
\emph{dukāndār har shām apnā hisāb likhtā hai.}
\end{quote}

\begin{quote}
Every evening the shopkeeper writes up his accounts.
\end{quote}

Masculine.

\begin{cousinweb}
\textbf{\ur{حساب}} (\emph{hisāb}) is usually traced, through Persian, to Arabic \emph{ḥisāb}, \textquotedblleft{}counting, reckoning\textquotedblright{}.
\end{cousinweb}

\begin{grammarlens}[title={What you've built}]
One more word for getting around.
\end{grammarlens}

\subsection*{Your turn}
\begin{itemize}
\raggedright
  \item \emph{Say it:} \emph{hisāb}
  \item \emph{Say it:} \emph{hisāb}, once more
  \item \emph{Say it:} \emph{dukāndār har shām apnā hisāb likhtā hai.} — Every evening the shopkeeper writes up his accounts.
  \item \emph{Recall:} say the Urdu for a cheque, then the Urdu for to withdraw (money), then say \emph{hisāb} again
\end{itemize}

\begin{tcolorbox}[breakable,colback=teal!4,colframe=teal!35!black,title={Before you move on}]
What does \ur{حساب} mean? (\textquotedblleft{}An account; a reckoning\textquotedblright{}.) Say it once more.
\end{tcolorbox}
\section[\texorpdfstring{\ur{سود} (sūd)}{سود (sūd)}]{\ur{سود} (sūd) — interest (on money)}

\label{lesson:UR-C418-sud}



\begin{quote}
Before the new one: say the Urdu for to withdraw (money), then the Urdu for a reckoning.
\end{quote}

\subsection*{You'll want to know: \ur{سود}}
\textbf{\ur{سود}} — \emph{sūd} — \textquotedblleft{}interest (on money)\textquotedblright{}.

\begin{quote}
\ur{بینک} \ur{قرض} \ur{پر} \ur{سود} \ur{لیتا} \ur{ہے۔}
\end{quote}

\begin{quote}
\emph{baink qarz par sūd letā hai.}
\end{quote}

\begin{quote}
The bank charges interest on a loan.
\end{quote}

Masculine.

\begin{cousinweb}
\textbf{\ur{سود}} (\emph{sūd}) is usually traced to Persian \emph{sūd}, \textquotedblleft{}profit, gain\textquotedblright{}; in Urdu it is the interest on a loan or on savings.
\end{cousinweb}

\begin{grammarlens}[title={What you've built}]
One more word for getting around.
\end{grammarlens}

\subsection*{Your turn}
\begin{itemize}
\raggedright
  \item \emph{Say it:} \emph{sūd}
  \item \emph{Say it:} \emph{sūd}, once more
  \item \emph{Say it:} \emph{baink qarz par sūd letā hai.} — The bank charges interest on a loan.
  \item \emph{Recall:} say the Urdu for to withdraw (money), then the Urdu for a reckoning, then say \emph{sūd} again
\end{itemize}

\begin{tcolorbox}[breakable,colback=teal!4,colframe=teal!35!black,title={Before you move on}]
What does \ur{سود} mean? (\textquotedblleft{}Interest (on money)\textquotedblright{}.) Say it once more.
\end{tcolorbox}
\section[\texorpdfstring{\ur{قسط} (qist)}{قسط (qist)}]{\ur{قسط} (qist) — an instalment}

\label{lesson:UR-C418-qist}



\begin{quote}
Before the new one: say the Urdu for a reckoning, then the Urdu for interest (on money).
\end{quote}

\subsection*{You'll want to know: \ur{قسط}}
\textbf{\ur{قسط}} — \emph{qist} — \textquotedblleft{}an instalment\textquotedblright{}.

\begin{quote}
\ur{میں} \ur{ہر} \ur{مہینے} \ur{موٹر} \ur{سائیکل} \ur{کی} \ur{قسط} \ur{دیتی} \ur{ہوں۔}
\end{quote}

\begin{quote}
\emph{maiṅ har mahīne moṭar sāikil kī qist detī hūṅ.}
\end{quote}

\begin{quote}
Every month I pay an instalment on the motorbike.
\end{quote}

Feminine.

\begin{cousinweb}
\textbf{\ur{قسط}} (\emph{qist}) is usually traced, through Persian, to Arabic \emph{qisṭ}, \textquotedblleft{}a portion, a share\textquotedblright{}.
\end{cousinweb}

\begin{grammarlens}[title={What you've built}]
One more word for getting around.
\end{grammarlens}

\subsection*{Your turn}
\begin{itemize}
\raggedright
  \item \emph{Say it:} \emph{qist}
  \item \emph{Say it:} \emph{qist}, once more
  \item \emph{Say it:} \emph{maiṅ har mahīne moṭar sāikil kī qist detī hūṅ.} — Every month I pay an instalment on the motorbike.
  \item \emph{Recall:} say the Urdu for a reckoning, then the Urdu for interest (on money), then say \emph{qist} again
\end{itemize}

\begin{tcolorbox}[breakable,colback=teal!4,colframe=teal!35!black,title={Before you move on}]
What does \ur{قسط} mean? (\textquotedblleft{}An instalment\textquotedblright{}.) Say it once more.
\end{tcolorbox}
\section[\texorpdfstring{\ur{ادائیگی} (adāigī)}{ادائیگی (adāigī)}]{\ur{ادائیگی} (adāigī) — a payment}

\label{lesson:UR-C418-adaigi}



\begin{quote}
Before the new one: say the Urdu for interest (on money), then the Urdu for an instalment.
\end{quote}

\subsection*{You'll want to know: \ur{ادائیگی}}
\textbf{\ur{ادائیگی}} — \emph{adāigī} — \textquotedblleft{}a payment\textquotedblright{}.

\begin{quote}
\ur{قرض} \ur{کی} \ur{ادائیگی} \ur{اگلے} \ur{سال} \ur{تک} \ur{ہے۔}
\end{quote}

\begin{quote}
\emph{qarz kī adāigī agle sāl tak hai.}
\end{quote}

\begin{quote}
The loan has to be paid back by next year.
\end{quote}

Feminine.

\begin{cousinweb}
\textbf{\ur{ادائیگی}} (\emph{adāigī}) is built on \emph{adā}, \textquotedblleft{}a payment, a carrying out\textquotedblright{}, usually traced, through Persian, to Arabic \emph{adāʾ}, \textquotedblleft{}discharging (a debt)\textquotedblright{}, with the Persian ending \emph{-gī}.
\end{cousinweb}

\begin{grammarlens}[title={What you've built}]
One more word for getting around.
\end{grammarlens}

\subsection*{Your turn}
\begin{itemize}
\raggedright
  \item \emph{Say it:} \emph{adāigī}
  \item \emph{Say it:} \emph{adāigī}, once more
  \item \emph{Say it:} \emph{qarz kī adāigī agle sāl tak hai.} — The loan has to be paid back by next year.
  \item \emph{Recall:} say the Urdu for interest (on money), then the Urdu for an instalment, then say \emph{adāigī} again
\end{itemize}

\begin{tcolorbox}[breakable,colback=teal!4,colframe=teal!35!black,title={Before you move on}]
What does \ur{ادائیگی} mean? (\textquotedblleft{}A payment\textquotedblright{}.) Say it once more.
\end{tcolorbox}
\section[\texorpdfstring{\ur{گھٹانا} (ghaṭānā)}{گھٹانا (ghaṭānā)}]{\ur{گھٹانا} (ghaṭānā) — to bring down (a price); to reduce}

\label{lesson:UR-C418-ghatana}



\begin{quote}
Before the new one: say the Urdu for an instalment, then the Urdu for a payment.
\end{quote}

\subsection*{You'll want to know: \ur{گھٹانا}}
\textbf{\ur{گھٹانا}} — \emph{ghaṭānā} — \textquotedblleft{}to bring down (a price); to reduce\textquotedblright{}.

\begin{quote}
\ur{بینک} \ur{نے} \ur{سود} \ur{گھٹایا۔}
\end{quote}

\begin{quote}
\emph{baink ne sūd ghaṭāyā.}
\end{quote}

\begin{quote}
The bank lowered its interest rate.
\end{quote}

Past, with \ur{نے} (he or she …): \ur{اس} \ur{نے} \ur{قیمت} \ur{گھٹائی}, \emph{us ne qīmat ghaṭāī}.

\begin{cousinweb}
\textbf{\ur{گھٹانا}} (\emph{ghaṭānā}) is built on \emph{ghaṭnā}, \textquotedblleft{}to decrease\textquotedblright{}, which CDIAL lists under \emph{ghaṭṭati}, a form it reconstructs (CDIAL 4415): this is the form for making something smaller.
\end{cousinweb}

\begin{grammarlens}[title={What you've built}]
One more word for getting around.
\end{grammarlens}

\subsection*{Your turn}
\begin{itemize}
\raggedright
  \item \emph{Say it:} \emph{ghaṭānā}
  \item \emph{Say it:} \emph{ghaṭānā}, once more
  \item \emph{Say it:} \emph{baink ne sūd ghaṭāyā.} — The bank lowered its interest rate.
  \item \emph{Recall:} say the Urdu for an instalment, then the Urdu for a payment, then say \emph{ghaṭānā} again
\end{itemize}

\begin{tcolorbox}[breakable,colback=teal!4,colframe=teal!35!black,title={Before you move on}]
What does \ur{گھٹانا} mean? (\textquotedblleft{}To bring down (a price); to reduce\textquotedblright{}.) Say it once more.
\end{tcolorbox}
